\section{A strict finite Fourier inequality with constant one}
\label{phase:section}

The next estimate improves the sequence exponent in the usual finite
Fourier bound while keeping the operator norm at most one. Keeping this
constant prevents a fixed loss from accumulating over the levels of the
angular hierarchy.
The distinction between a dominant coefficient and a strict norm
deficit is related to the study of near-extremizers for discrete Young
and Hausdorff--Young inequalities. Relevant precedents include
Fournier's work on sharpness, the fourth-moment case of Eisner and Tao,
and the quantitative discrete results of Charalambides and Christ
\cite{Fournier1977,EisnerTao2012,CharalambidesChrist2011}.
Here the frequencies are arbitrary separated real numbers and the
averaging probability is specified by its Fourier support. We prove
the needed quantitative gap in that setting, including the exact
constant one after the coefficient exponent is improved.

Throughout this section the Fourier transform of a finite measure is
\(\widehat\nu(\omega)=\int e^{-i\omega t}\,\dd\nu(t)\).

Fix the absolute constants
\begin{equation}
 \epsilon_0=\frac14,
 \qquad \kappa_0=2^{-10},
 \qquad c_{\rm ph}=2^{-60}.
 \label{phase:absolute-constants}
\end{equation}
For the bound \(D\ge1\) from the preceding section and for \(p\ge4\), put
\begin{equation}
 \alpha=2+\frac{c_{\rm ph}}{\log(2D)},
 \qquad \mu=\frac{\alpha-1}{p},
 \qquad r_* = \frac{1}{1-\mu}.
 \label{phase:parameters}
\end{equation}
In particular, \(2<\alpha<7/3\), \(0<\mu<1/3\), and
\(1<r_*\le3/2\). The numerical values in
\eqref{phase:absolute-constants} are convenient fixed choices satisfying
the estimates below.

\begin{proposition}[Finite Fourier gain]
\label{phase:finite-fourier}
Let \(\lambda_1,\ldots,\lambda_d\in\R\), where \(d\le D\), satisfy
\(\abs{\lambda_i-\lambda_j}\ge\Delta>0\) for \(i\ne j\).
Let \(\nu\) be any positive probability measure on \(\R\) such that
\(\supp\widehat\nu\subset(-\epsilon_0\Delta,\epsilon_0\Delta)\).
Then, for every finite \(p\ge4\) and every complex coefficient sequence,
\begin{equation}
 \left(\int_\R
       \left|\sum_{i=1}^d a_i e^{i\lambda_i t}\right|^p
       \,\dd\nu(t)\right)^{1/p}
 \le \left(\sum_{i=1}^d |a_i|^{r_*}\right)^{1/r_*}.
 \label{eq:src-13}
\end{equation}
The same estimate holds for coefficients in a finite-dimensional Hilbert
space \(\mathcal H\) when the exponentials act diagonally: in an
orthonormal basis, the \(\xi\)-th component may use frequencies
\(\lambda_{i,\xi}\), provided they are \(\Delta\)-separated as \(i\)
varies, for each fixed \(\xi\).  The right side is then
\( (\sum_i\norm{a_i}_{\mathcal H}^{r_*})^{1/r_*}\).
All these estimates have constant one.
\end{proposition}

We first prove a quantitative gap in the fourth moment when no
coefficient dominates. The proof compares the law of three indices
subject to an approximate additive relation with the independent triple
law; separated lower and upper quantiles force a definite difference
between these laws. We then handle a dominant coefficient directly,
interpolate, and pass to Hilbert-valued coefficients.

\begin{lemma}[A quantitative gap away from a single coefficient]
\label{phase:quantile-gap}
For a separated frequency list as above, define
\[
 \mathcal Q=\{(i,j,k,l):
    |\lambda_i+\lambda_j-\lambda_k-\lambda_l|
       <\epsilon_0\Delta\}.
\]
Let \(x_i\ge0\), \(\sum_i x_i=1\), and put
\[
 Q(x)=\sum_{(i,j,k,l)\in\mathcal Q}
                (x_i x_j x_k x_l)^{3/4}.
\]
Then \(Q(x)\le1\).  If \(0<\kappa<1\) and
\(\max_i x_i\le1-\kappa\), then
\begin{equation}
 Q(x)\le1-\delta_\kappa,
 \qquad
 \delta_\kappa=
       \frac{\kappa^4(1-\kappa/2)^2}{4096}
       \ge 2^{-14}\kappa^4.
 \label{phase:quantitative-gap}
\end{equation}
\end{lemma}

\begin{proof}
Any three specified indices have at most one completion to an element of
\(\mathcal Q\): two possible missing frequencies would be within
\(2\epsilon_0\Delta<\Delta\) of each other.  On the finite set of
quadruples, consider the four nonnegative measures
\[
 A=x_i x_j x_k\one_{\mathcal Q},\quad
 B=x_i x_k x_l\one_{\mathcal Q},\quad
 C=x_i x_j x_l\one_{\mathcal Q},\quad
 E=x_j x_k x_l\one_{\mathcal Q}.
\]
Each has total mass at most one, by unique completion.  Their pointwise
geometric mean is the summand defining \(Q(x)\), so H\"older's inequality
already gives \(Q(x)\le1\).

For nonnegative measures on a finite set write
\(\mathsf H(f,g)=\sum\sqrt{fg}\) for their Hellinger affinity.
Cauchy--Schwarz gives
\begin{equation}
 Q(x)\le
       \bigl(\mathsf H(A,B)\mathsf H(C,E)\bigr)^{1/2}
       \le \mathsf H(A,B)^{1/2}.
 \label{phase:affinity-lower-bound}
\end{equation}
Project \(A\) to the three coordinates \((i,k,l)\), obtaining a measure
\(P\) of mass at most one, and let
\(P_0(i,k,l)=x_i x_k x_l\) be the independent triple law.
For each \((i,k,l)\) there is at most one completing \(j\).  Consequently
\begin{equation}
 \mathsf H(P,P_0)=\mathsf H(A,B)\ge Q(x)^2.
 \label{phase:projected-affinity}
\end{equation}
This identity also holds when a factor \(x_i\), \(x_k\), or \(x_l\)
vanishes, since the corresponding terms on both sides are zero.

Adjoin one point to the triple space and place the missing mass of \(P\)
there, giving a probability \(P^+\).  Extend \(P_0\) by zero at that
point.  For any two probabilities \(U,V\),
\[
 \frac12\sum |U-V|
 \le \frac12
       \left(\sum(\sqrt U-\sqrt V)^2\right)^{1/2}
       \left(\sum(\sqrt U+\sqrt V)^2\right)^{1/2}
 =\sqrt{1-\mathsf H(U,V)^2}.
\]
Every event \(\mathcal A\) in the triple space therefore satisfies
\begin{equation}
 P_0(\mathcal A)-P(\mathcal A)
 \le \sqrt{1-\mathsf H(P,P_0)^2}
 \le \sqrt{1-Q(x)^4}.
 \label{phase:event-deficit-upper}
\end{equation}
Thus a fourth moment close to the ordinary bound would force the
completed triple law close to the independent law, quantitatively and
without a factor depending on the number of frequencies.

Let \(X\) be a random frequency with probabilities \(x_i\).
Choose lower and upper \(\kappa/4\)-quantile frequencies \(a,b\), so that
\[
 \Pr(X\le a),\Pr(X\ge b)\ge\kappa/4,
 \qquad
 \Pr(X<a),\Pr(X>b)\le\kappa/4.
\]
They can be chosen with \(a\le b\).  If \(a=b\), that atom would have
mass at least \(1-\kappa/2>1-\kappa\), a contradiction.  Hence
\(b-a\ge\Delta\), and \([a,b]\) has mass at least \(1-\kappa/2\).
Of its two closed halves, choose one, denoted \([e,h]\), with mass at
least \((1-\kappa/2)/2\).

Consider the triple event
\[
 \mathcal A=\{\lambda_i\ge b,\ \lambda_k\le a,\ \lambda_l\le h\}.
\]
For any completion in \(\mathcal Q\) with these pair restrictions,
\[
 \lambda_l-\lambda_j
 >\lambda_i-\lambda_k-\epsilon_0\Delta
 \ge b-a-\epsilon_0\Delta
 > (b-a)/2=h-e.
\]
Thus \(\lambda_l\le h\) forces \(\lambda_j<e\).  Under the measure
\(A\), the free coordinates \((i,j,k)\) carry their independent
probabilities, with a possible deletion when no completion exists.
It follows that
\[
 P(\mathcal A)
 \le \Pr(X\ge b)\Pr(X\le a)\Pr(X<e),
\]
whereas
\[
 P_0(\mathcal A)
 =\Pr(X\ge b)\Pr(X\le a)\Pr(X\le h).
\]
Their difference is at least
\[
 g_\kappa=
       (\kappa/4)^2\frac{1-\kappa/2}{2}
       =\frac{\kappa^2(1-\kappa/2)}{32}.
\]
Combining this with \eqref{phase:event-deficit-upper} gives
\(Q(x)^4\le1-g_\kappa^2\).  Since
\((1-t)^{1/4}\le1-t/4\) for \(0\le t\le1\),
\(Q(x)\le1-g_\kappa^2/4\).  This is
\eqref{phase:quantitative-gap}.
\end{proof}

\begin{proof}[Proof of \Cref{phase:finite-fourier}]
\emph{The ordinary fourth-moment bound.}
Write \(Ta(t)=\sum_i a_i e^{i\lambda_i t}\) and \(s=4/3\).
The Fourier support hypothesis and positivity of \(\nu\) give
\[
 \int |Ta|^4\,\dd\nu
 \le \sum_{\mathcal Q}|a_i a_j a_k a_l|.
\]
Indeed every other term in the fourth-power expansion averages to zero,
and the absolute value of every surviving Fourier coefficient is at most
one.  If \(\norm a_{\ell^s}=1\), the last sum is \(Q(x)\) for
\(x_i=|a_i|^s\).  Thus \Cref{phase:quantile-gap} proves the ordinary
norm-one bound
\begin{equation}
 \norm{Ta}_{L^4(\nu)}\le\norm a_{\ell^{4/3}}.
 \label{phase:ordinary-fourth}
\end{equation}

For the improved bound at four, put
\[
 \theta=\frac{c_{\rm ph}}{\log(2D)},
 \qquad r_4=\frac4{3-\theta}.
\]
Then \(4/3<r_4\le3/2\) and
\begin{equation}
 \norm a_{\ell^{4/3}}
 \le D^{\theta/4}\norm a_{\ell^{r_4}},
 \qquad D^\theta\le e^{c_{\rm ph}}.
 \label{phase:exponent-comparison}
\end{equation}
\emph{No dominant coefficient.}
For a nonzero sequence, normalize \(x_i=|a_i|^{4/3}/
\sum_k|a_k|^{4/3}\).  If \(\max_i x_i\le1-\kappa_0\),
\eqref{phase:quantitative-gap} and
\eqref{phase:exponent-comparison} yield
\[
 \norm{Ta}_{L^4(\nu)}^4
 \le (1-\delta_{\kappa_0})e^{c_{\rm ph}}
                 \norm a_{\ell^{r_4}}^4
 \le \norm a_{\ell^{r_4}}^4.
\]
Here \(\delta_{\kappa_0}\ge2^{-54}>c_{\rm ph}\), and
\(1-\delta\le e^{-\delta}\).

\emph{A dominant coefficient.}
Suppose instead that one coefficient has normalized mass exceeding
\(1-\kappa_0\).  Divide all coefficients by that coefficient and remove
its carrier from \(Ta\).  The normalized sum is \(1+U(t)\), and, if
\(b\) denotes the \(\ell^{r_4}\) norm of the other coefficients,
\begin{equation}
 b\le\left(\frac{\kappa_0}{1-\kappa_0}\right)^{3/4}<\frac1{64}.
 \label{phase:small-remainder}
\end{equation}
Frequency separation gives the exact identities
\[
 \int U\,\dd\nu=0,
 \qquad
 \int |U|^2\,\dd\nu
      =\sum_{i\ne i_0}|a_i/a_{i_0}|^2\le b^2.
\]
By \eqref{phase:ordinary-fourth} and
\eqref{phase:exponent-comparison},
\(\int|U|^4\,\dd\nu\le e^{c_{\rm ph}}b^4\), and hence
\(\int|U|^3\,\dd\nu\le e^{c_{\rm ph}/2}b^3\).
Expanding \(|1+U|^4\), the linear term has mean zero.  Bounding the
remaining terms gives
\begin{align*}
 \int|1+U|^4\,\dd\nu
 &\le1+6b^2+4e^{c_{\rm ph}/2}b^3+e^{c_{\rm ph}}b^4\\
 &\le1+14b^2
 \le1+\frac83 b^{r_4}
 \le(1+b^{r_4})^{4/r_4}.
\end{align*}
For the penultimate inequality, use
\(b^{2-r_4}\le\sqrt b\le1/8\); for the last one, use
\(4/r_4\ge8/3\) and convexity.  Restoring the dominant coefficient
proves the norm-one \(\ell^{r_4}\)-to-\(L^4(\nu)\) estimate in this
case too.  The zero sequence is immediate.

\emph{Interpolation with constant one.} The same
operator has norm at most one from \(\ell^1\) to \(L^\infty(\nu)\).
For \(0<t\le1\), interpolation between these two bounds gives
\[
 \frac1{p_t}=\frac t4,
 \qquad
 \frac1{r_t}=1-t+\frac t{r_4}.
\]
One can see the normalization directly from the three-lines lemma
\cite[Lemma~3.3, p.~57]{Hajlasz2012}.
For \(\norm a_{\ell^{r_t}}=\norm g_{L^{p_t'}(\nu)}=1\), initially
with \(g\) simple, interpolate their nonzero magnitudes by
\[
 a_i(z)=\operatorname{sgn}(a_i)
          |a_i|^{r_t(1-z+z/r_4)},
 \qquad
 g(z)=\operatorname{sgn}(g)
          |g|^{p_t'(1-z+3z/4)}.
\]
Here \(\operatorname{sgn}\) denotes complex phase, and the phase of
\(g\) is chosen for a bilinear dual pairing.  The zero values stay zero.
On the two vertical boundaries, these functions have respective unit
norms in \((\ell^1,L^1)\) and \((\ell^{r_4},L^{4/3})\).
The analytic function \(\int (Ta(z))g(z)\,\dd\nu\) is thus bounded by
one there and at \(z=t\).  Finite sums and simple functions justify
analyticity and boundedness on the strip; density and duality remove
the simplicity restriction.  Taking \(t=4/p\) gives
\[
 \frac1{r_t}
 =1-\frac{1+\theta}{p}
 =1-\frac{\alpha-1}{p},
\]
which is precisely \(r_t=r_*\).  This proves the scalar form of
\eqref{eq:src-13} for every finite \(p\ge4\).

\emph{Hilbert-valued coefficients.}
Finally write the Hilbert space as \(\ell^2\) in an orthonormal basis,
and let \(F_\xi(t)=\sum_i a_{i,\xi}e^{i\lambda_{i,\xi}t}\).
Minkowski's inequality, the scalar result, and Minkowski's inequality
again give
\begin{align*}
 \norm F_{L^p(\nu;\ell^2)}
 &\le\left(\sum_\xi\norm{F_\xi}_{L^p(\nu)}^2\right)^{1/2}\\
 &\le\left(\sum_\xi
          \left(\sum_i|a_{i,\xi}|^{r_*}\right)^{2/r_*}
            \right)^{1/2}\\
 &\le\left(\sum_i
          \left(\sum_\xi|a_{i,\xi}|^2\right)^{r_*/2}
            \right)^{1/r_*}.
\end{align*}
The first inequality uses \(p\ge2\), and the last uses \(r_*\le2\).
Each has constant one and is independent of the Hilbert dimension.
Only the frequency list within an individual component enters the
scalar estimate, completing the proof.
\end{proof}
